\documentclass[11pt]{article}
\usepackage{fontspec}
\usepackage[french]{babel}
\usepackage{amsmath}
\usepackage{amssymb}
\usepackage{amsthm}
\usepackage{mathrsfs}
\usepackage{tikz-cd}
\usepackage{enumitem}
\usetikzlibrary{babel}

\newcommand{\printedpage}[1]{\marginpar{\footnotesize p.~#1}}
\newcommand{\msp}[1]{\marginpar{\footnotesize [#1]}}

\newcommand{\HOT}{\mathrm{HOT}}
\newcommand{\Cat}{\mathrm{Cat}}
\newcommand{\Hot}{\mathrm{Hot}}
\newcommand{\hot}{\mathrm{hot}}
\newcommand{\Ob}{\operatorname{Ob}}
\newcommand{\Hom}{\operatorname{Hom}}
\newcommand{\ul}[1]{\underline{\mathrm{#1}}}
\newcommand{\cercle}[1]{\tikz[baseline=(c.base)]\node[draw,circle,inner sep=1pt](c){#1};}
\newcommand{\ovale}[1]{\tikz[baseline=(c.base)]\node[draw,rounded corners=0.8ex,inner sep=2pt](c){#1};}

\begin{document}
\setcounter{footnote}{45}

\printedpage{44}
% head: A. Grothendieck : Dérivateurs, Ch. VI, Édition des Universités de Montpellier II et Paris VII
% foot: Version 9 novembre 2001
et les localiseurs $W_S^{\mathrm{g}}$ dans ces catégories ($f\in W_S^{\mathrm{g}}$ si et seulement si les $f/s : X/s\longrightarrow Y/s$ sont dans $W$). On trouve une `même catégorie localisée', notée
\[ \HOT^o_W(S), \]
et on a tautologiquement
\msp{page 71}
\[ \HOT^o_W(S)\approx\HOT_W(S^o). \]
\texttt{[plutôt} $(\HOT_W(S^o))^o$ \texttt{cf. p. 50]}

\ovale{III}\quad Sur les six catégories internes du diagramme, les trois notions de $W$-équivalence, en termes de $W_S$, $W_S^{\mathrm{d}}$, $W_S^{\mathrm{g}}$ coïncident, et coïncident avec la $W$-équivalence par fibres $W_S^{\mathrm{f}}$. Sur les trois catégories de $(*)$ autres que $\Cat/S$, les localiseurs $W_S^{\mathrm{f}}$, $W_S^{\mathrm{d}}$ coïncident; sur les trois catégories de $(*')$ autres que $\Cat/S$, les localiseurs $W_S^{\mathrm{f}}$, $W_S^{\mathrm{g}}$ coïncident. Mais en général (\footnote{sans doute dans tous les cas où $S$ n'est pas équivalente à une catégorie dicrète.% [sic] dicrète -> discrète, faute d'orthographe
}), sur ces six catégories, ces localiseurs $W_S^{\mathrm{f}}$ ne sont pas égaux à $W_S$, plus précisément donnent des localisées non équivalentes, elles sont donc munies chacune de \textit{deux} localiseurs essentiellement distincts, $W_S^{\mathrm{f}}$ et $W_S$. La catégorie $\Cat/S$ est la seule catégorie du diagramme sur laquelle les localiseurs $W_S$, $W_S^{\mathrm{d}}$ et $W_S^{\mathrm{g}}$ sont tous les trois essentiellment distincts. % [sic] essentiellment -> essentiellement, lettre omise
C'est la seule qui permette d'obtenir simultanément les trois catégories localisées, comme
\[
\begin{tikzcd}[column sep=small]
 & \HOT^{\mathrm{lc}}_W(S)\simeq(\Cat/S)W_S^{-1} & \\
\HOT_W(S)\simeq(\Cat/S)(W_S^{\mathrm{d}})^{-1} \arrow[ur,"\rho"] & & \HOT^o_W(S)\simeq(\Cat/S)(W_S^{\mathrm{g}})^{-1} \arrow[ul,"\sigma"']
\end{tikzcd}
\]

Ici, $\rho$ et $\sigma$ sont des foncteurs de localisation (comme il est évident sur la définition).

Tout foncteur ou foncteur composé du diagramme page 8, allant d'une catégorie interne au diagramme vers une catégorie bordante, induit un foncteur pour les localisées pour $W_S^{\mathrm{d}}$ (quand la catégorie but fait partie des quatre gauches,
\msp{page 72}
rappelées dans $(*)$, p. 70), soit pour $W_S^{\mathrm{g}}$ (quand ces catégories font partie des quatre du côté droit, rappelées dans $(*')$), lesquels localiseurs coïncident aussi avec $W_S^{\mathrm{f}}$ sauf dans le seul cas de la catégorie $\Cat/S$. Tous les foncteurs ainsi obtenus se réduisent, à équivalence près, à deux, suivant que la catégorie localisée but est $\HOT_W(S)$ (cas de la localisation $W_S^{\mathrm{d}}$), ou qu'elle est $\HOT^o_W(S)$ (cas de la localisation $W_S^{\mathrm{g}}$). On trouve ainsi des foncteurs (en sens inverse de $\rho$ et de $\sigma$)
\printedpage{45}
% head: A. Grothendieck : Dérivateurs, Ch. VI, Édition des Universités de Montpellier II et Paris VII
% foot: Version 9 novembre 2001
\[
\begin{tikzcd}
 & \HOT^{\mathrm{lc}}_W(S) \arrow[dl,"i"'] \arrow[dr,"j"] & \\
\HOT_W(S) & & \HOT^o_W(S).
\end{tikzcd}
\]

Les foncteurs $\rho$, $\sigma$ s'interchangent par dualité $S\longmapsto S^o$, et de même pour $i$, $j$.

\ovale{IV}\quad Les foncteurs $\rho$, $i$ sont adjoints l'un de l'autre dans cet ordre, i.e. on a pour deux objets $\xi\in\Ob\HOT_W(S)$, $\eta\in\Ob\HOT^{\mathrm{lc}}_W(S)$
\[ \Hom(\rho(\xi),\eta)\simeq\Hom(\xi,i(\eta)) \]
(isomorphisme de bifoncteurs). Et dualement les foncteurs $\sigma$, $j$ sont adjoint l'un de l'autre % [sic] adjoint -> adjoints, accord du pluriel
dans cet ordre : pour $\xi\in\Ob\HOT^o_W(S)$, $\eta\in\Ob\HOT^{\mathrm{lc}}_W(S)$
\[ \Hom(\xi,j(\eta))\simeq\Hom(\sigma(\xi),\eta). \]

\msp{page 73}
Il s'ensuit notamment que \textit{les foncteurs $i$, $j$ sont pleinement fidèles.}

\cercle{V}\quad On a des foncteurs évidents
\[
\begin{aligned}
&(*) &\qquad& \HOT_W(S) &\quad& \longrightarrow &\quad& \underline{\mathrm{Hom}}(S^o,\Hot_W)\\
&(*') && \HOT^o_W(S) && \longrightarrow && \underline{\mathrm{Hom}}(S,\Hot_W).
\end{aligned}
\]
Le premier s'interprète sur les quatre catégories $(*)$ (bordantes gauches) de page 70, qui servent indifférement à décrire $\HOT_W(S)$, par % [sic] indifférement -> indifféremment, faute d'orthographe
\[ s\longmapsto\hot_W(X_s) \]
dans le cas des trois catégories distinctes de $\Cat/S$, et par
\[ s\longmapsto\hot_W(s\backslash X) \]
dans le cas de $\Cat/S$. La dépendence fonctorielle par rapport à $s$ est évidente dans le cas % [sic] dépendence -> dépendance, faute d'orthographe
de $\Cat/S$ (donnée par $s\longmapsto\hot_W(s\backslash X)$), et dans le cas de $\ul{Fibsc}\,S$, $\ul{Fib}\,S$, où on dispose, pour $u : s\longrightarrow t$ dans $S$, du changement de base
\[ u^* : X_t\longrightarrow X_s. \]
Dans le cas de $\ul{Liss}\,S$, la dépendence fonctorielle se déduit de celle de $\hot_W(s\backslash X)$, grâce % [sic] dépendence -> dépendance, faute d'orthographe
au fait que l'inclusion
\[ X_s\longrightarrow s\backslash X \]
est alors une $W$-équivalence.

On a une interprétation duale de l'application
\printedpage{46}
% head: A. Grothendieck : Dérivateurs, Ch. VI, Édition des Universités de Montpellier II et Paris VII
% foot: Version 9 novembre 2001
\msp{page 74}
$(*)'$, \texttt{[mot illisible]}
\[ s\longmapsto\hot_W(X_s)\qquad\text{ou}\qquad s\longmapsto\hot_W(s\backslash X), \]
avec la loi fonctorielle en $s$ déduite du foncteur $t\backslash X\longrightarrow s\backslash X$ \texttt{[plutôt} $X/s\longrightarrow X/t$\texttt{]} (pour $u : s\longrightarrow t$) dans le cas du $\Cat/S$, du foncteur cochangement de base dans les cas de $\ul{Cofib}\,S$ et $\ul{Cofibsc}\,S$, et de la variante pour $\Cat/S$ et de la $W$-équivalence
\[ X_s\longrightarrow X/s \]
dans le cas de $\ul{Prop}\,S$.

On trouve alors ceci : pour qu'un objet $\xi$ de $\HOT_W(S)$ (resp. $\HOT^o_W(S)$) soit dans l'image essentielle de $\HOT^{\mathrm{lc}}_W(S)$, il faut et il suffit que le foncteur associé $S^o\longrightarrow\Hot_W$ (resp. $S\longrightarrow\Hot_W$) soit un système local (\footnote{il faudra encore écrire la vérification}), i.e. transforme toute flèche de $S$ en un isomorphisme, i.e. se factorise par $(\Pi S)^o\longrightarrow\Hot_W$ resp. $\Pi S\longrightarrow\Hot_W$ (\footnote{où $\Pi S$ est le groupoïde fondamental de $S$}). Donc pour un $X$ dans $\Cat/S$, cela signifie soit que les foncteurs
\[ t\backslash X\longrightarrow s\backslash X \]
associés à des $u : s\longrightarrow t$ dans $S$ soient tous des $W$-équivalences (cas de $\HOT_W(S)$), soit que les foncteurs
\[ X/s\longrightarrow X/t \]
le soient (cas de $\HOT^o_W(S)$). (NB. ces conditions ne sont nullement équivalentes. Mais si $X$ \texttt{[est]} lisse \texttt{[sur]} $S$,
\msp{page 75}
la première implique la seconde, et l'inverse est vrai si $X$ est propre sur $S$, et dans ce cas, la condition équivaut à ce que $X$ soit dans $\ul{Fibhot}$. Enfin, si $X$ est déjà dans $\ul{Fibhot}\,S$, alors les deux conditions sont automatiquement satisfaites (\footnote{On aurait dû dans le diagramme introduire aussi $\ul{Liss}_0S$, $\ul{Prop}_0S$ !}).
\[
\begin{tikzcd}[column sep=small, row sep=small]
\ul{Fibsc}_0S \arrow[dd,hook] \arrow[drr] & & & & & & & & \ul{Cofibsc}_0S \arrow[dll] \arrow[dd,hook] \\
 & & \ul{Fib}_0S \arrow[dd,hook] \arrow[dr,hook] & & \ul{Parf}\,S \arrow[dl,hook] \arrow[dd,hook] \arrow[dr,hook] & & \ul{Cofib}_0S \arrow[dl,hook] \arrow[dd,hook] & & \\
\ul{Fibsc}\,S \arrow[drr] & & & \ul{Liss}_0S \arrow[dr,hook] \arrow[dd,hook] & & \ul{Prop}_0S \arrow[dl,hook] \arrow[dd,hook] & & & \ul{Cofibsc}\,S \arrow[dll] \\
 & & \ul{Fib}\,S \arrow[dr,hook] & & \ul{Fibhot}\,S \arrow[dd,hook] \arrow[uuullll,dotted,bend right,"\Phi_0"'] \arrow[uuurrrr,dotted,bend left,"\Psi_0"'] & & \ul{Cofib}\,S \arrow[dl,hook] & & \\
 & & & \ul{Liss}\,S \arrow[dr,hook] & & \ul{Prop}\,S \arrow[dl,hook] & & & \\
 & & & & \Cat/S \arrow[uuullll,dotted,bend left,"\Phi"] \arrow[uuurrrr,dotted,bend right,"\Psi"'] & & & &
\end{tikzcd}
\]
\printedpage{47}
% head: A. Grothendieck : Dérivateurs, Ch. VI, Édition des Universités de Montpellier II et Paris VII
% foot: Version 9 novembre 2001
(\footnote{Le diagramme de \textit{quinze} catégories remarquables.}). De plus, si deux éléments $\xi\in\Ob\HOT_W(S)$, $\xi'\in\Ob\HOT^o_W(S)$ proviennent d'un même $\eta\in\Ob\HOT^{\mathrm{lc}}_W(S)$, alors les deux foncteurs correspondants
\[ (\Pi S)^o\xrightarrow{\ \varphi\ }\Hot_W\xleftarrow{\ \varphi'\ }\Pi S \]
sont \textit{associés,} i.e. chacun est isomorphe à `l'inverse' de l'autre.
\[
\left\{
\begin{array}{lcl}
\varphi^o(s) & = & \varphi(s)\\
\varphi^o(u) & = & \varphi(u)^{-1}.
\end{array}
\right.
\]

\ovale{VI}\quad NB Dans la nouvelle version du diagramme, où j'ai rajouté des catégories $\ul{Liss}_0S$ et $\ul{Prop}_0S$ pour rendre le tracé plus transparent, les catégories appelés `bordantes' deviennent % [sic] appelés -> appelées, accord du participe
le `rez-de-chausseée', et les catégories `internes', le `premier étage' ou `étage'. Tous % [sic] rez-de-chausseée -> rez-de-chaussée, lettre en trop
les 7 carrés du diagramme sont catésiens. Toutes les flèches sont des inclusions pleinement % [sic] catésiens -> cartésiens, lettre omise
fidèles, sauf les quatre flèches horizontales extérieures des catégories de Cat-fibrés ou de Cat-cofibrés \textit{scindés}, avec morphismes exclusivement cartésiens ou cocartésiens et même compatibles avec les scindages, vers les catégories correspondantes de structures non scindées.

\msp{page 76}
\[
\begin{tikzcd}[column sep=small, row sep=small]
 & & & & \ul{Fibhot}\,S \arrow[dd,hook,"\delta"] \arrow[dddllll,dotted,bend right,"\Phi_0"'] \arrow[dddrrrr,dotted,bend left,"\Psi_0"] & & & & \\
 & & & \ul{Liss}_0S \arrow[ur,hook] \arrow[dd,hook] & & \ul{Prop}_0S \arrow[ul,hook] \arrow[dd,hook] & & & \\
 & \ul{Fib}_0S \arrow[urr,hook] \arrow[dd,hook] \arrow[ddrrr,dotted,bend left] & & & \Cat/S \arrow[dd,"C"] \arrow[dddllll,dotted,bend right,"\Phi"'] \arrow[dddrrrr,dotted,bend left,"\Psi"] & & & \ul{Cofib}_0S \arrow[ull,hook] \arrow[dd,hook] \arrow[ddlll,dotted,bend right] & \\ % Reading verified against PDF leaf 9: sources Fib_0 S and Cofib_0 S, target middle Parf S.
\ul{Fibsc}_0S \arrow[ur] \arrow[dd,hook] & & & \ul{Liss}\,S \arrow[ur,hook] & & \ul{Prop}\,S \arrow[ul,hook] & & & \ul{Cofibsc}_0S \arrow[ul] \arrow[dd,hook] \\
 & \ul{Fib}\,S \arrow[urr,hook] & & & \ul{Parf}\,S \arrow[uuul] \arrow[uuur] \arrow[dd,dotted,"\mathrm{id}"] & & & \ul{Cofib}\,S \arrow[ull,hook] & \\
\ul{Fibsc}\,S \arrow[ur] & & & & & & & & \ul{Cofibsc}\,S \arrow[ul] \\
 & & & & \ul{Parf}\,S \arrow[uuul,dotted,hook] \arrow[uuur,dotted,hook] & & & &
\end{tikzcd}
\]

J'ai tracé en pointillés verts \texttt{[couleur dans l'original seulement]} les cinq foncteurs $\Phi$, $\Psi$, $\Phi_0$, $\Psi_0$, $C$ qui par passage aux localisés donnent des quasi-inverses pour toutes les équivalences de catégories énoncées dans \cercle{I} \cercle{II} du présent Scholie. De façon précise, le foncteur
\[ \Phi : \Cat/S\longrightarrow\ul{Fibsc}\,S \]
associe à toute $X$ dans $\Cat/S$ la catégorie scindée sur $S$
\[ s\longmapsto s\backslash X \]
\printedpage{48}
% head: A. Grothendieck : Dérivateurs, Ch. VI, Édition des Universités de Montpellier II et Paris VII
% foot: Version 9 novembre 2001
sur $S$. Ce foncteur, et ses restrictions à $\ul{Liss}\,S$, $\ul{Fib}\,S$ et, plus généralement, à toute sous-catégorie pleine $\mathcal{F}$ de $\Cat/S$ contenant $\ul{Fib}\,S$, est compatible avec les localiseurs $W_S$ d'une part (sur $\mathcal{F}$ et $\ul{Fibsc}\,S$), $W_S^{\mathrm{d}}$ d'autre part, et passe donc aux localisés correspondants, pour fournir le quasi-inverse cherché de
\[ (\ul{Fibsc}\,S)\Sigma^{-1}\longrightarrow\mathcal{F}\Sigma^{-1}\qquad(\Sigma=W_S\text{ ou }W_S^{\mathrm{d}}) \]

On peut aussi, pour toute sous-catégorie pleine $\mathcal{F}'$ intermédiaire
\[ \ul{Fib}\,S\subseteq\mathcal{F}'\subseteq\mathcal{F}\subseteq\Cat/S, \]
regarder $\Phi$ comme un foncteur $\mathcal{F}\longrightarrow\mathcal{F}'$ (plus précisément, prendre le composé
\[ \mathcal{F}\xrightarrow{\Phi|\mathcal{F}}\ul{Fibsc}\,S\longrightarrow\ul{Fib}\,S\hookrightarrow\ \mathcal{F}'\ ), \]
soit $\Phi_{\mathcal{F}',\mathcal{F}}$, et de ce qui vient d'être dit résulte que les deux foncteurs
\[
\begin{tikzcd}
\mathcal{F}' \arrow[r,hook] & \mathcal{F} \arrow[l,bend left,"\Phi_{\mathcal{F}',\mathcal{F}}"]
\end{tikzcd}
\]
\msp{page 77}
établissent des équivalences quasi-inverses l'une de l'autre
\[
\begin{tikzcd}
\mathcal{F}'\Sigma^{-1} \arrow[r,"\approx"] & \mathcal{F}\Sigma^{-1} \arrow[l,bend left,"\approx"',"\Phi_{\mathcal{F}',\mathcal{F}}"]
\end{tikzcd}
\qquad\qquad(\Sigma=W_S\text{ ou }W_S^{\mathrm{d}})
\]

Énoncé dual pour le foncteur
\[ \Psi : \Cat/S\longrightarrow\ul{Cofibsc}\,S, \]
associant à $X$ dans $\Cat/S$ la catégorie cofibrée scindée sur $S$
\[ s\longmapsto X/s. \]
Il donne des couples de foncteurs quasi-inverses
\[
\begin{tikzcd}
(\ul{Cofibsc}\,S)\Sigma^{-1} \arrow[r] & \mathcal{F}\Sigma^{-1} \arrow[l,bend left,"\Psi_{\mathcal{F}}"]
\end{tikzcd}
\qquad\qquad(\Sigma=W_S\text{ ou }W_S^{\mathrm{g}})
\]
pour tout sous-catégorie pleine $\mathcal{F}$ de $\Cat/S$ telle que % [sic] tout -> toute, accord avec sous-catégorie
\[ \ul{Cofib}\,S\subseteq\mathcal{F}\subseteq\Cat/S, \]
et des couples de foncteurs quasi-inverses

\end{document}
